% Options for packages loaded elsewhere
\PassOptionsToPackage{unicode}{hyperref}
\PassOptionsToPackage{hyphens}{url}
\documentclass[
]{article}
\usepackage{xcolor}
\usepackage{amsmath,amssymb}
\setcounter{secnumdepth}{-\maxdimen} % remove section numbering
\usepackage{iftex}
\ifPDFTeX
  \usepackage[T1]{fontenc}
  \usepackage[utf8]{inputenc}
  \usepackage{textcomp} % provide euro and other symbols
\else % if luatex or xetex
  \usepackage{unicode-math} % this also loads fontspec
  \defaultfontfeatures{Scale=MatchLowercase}
  \defaultfontfeatures[\rmfamily]{Ligatures=TeX,Scale=1}
\fi
\usepackage{lmodern}
\ifPDFTeX\else
  % xetex/luatex font selection
\fi
% Use upquote if available, for straight quotes in verbatim environments
\IfFileExists{upquote.sty}{\usepackage{upquote}}{}
\IfFileExists{microtype.sty}{% use microtype if available
  \usepackage[]{microtype}
  \UseMicrotypeSet[protrusion]{basicmath} % disable protrusion for tt fonts
}{}
\makeatletter
\@ifundefined{KOMAClassName}{% if non-KOMA class
  \IfFileExists{parskip.sty}{%
    \usepackage{parskip}
  }{% else
    \setlength{\parindent}{0pt}
    \setlength{\parskip}{6pt plus 2pt minus 1pt}}
}{% if KOMA class
  \KOMAoptions{parskip=half}}
\makeatother
\usepackage{color}
\usepackage{fancyvrb}
\newcommand{\VerbBar}{|}
\newcommand{\VERB}{\Verb[commandchars=\\\{\}]}
\DefineVerbatimEnvironment{Highlighting}{Verbatim}{commandchars=\\\{\}}
% Add ',fontsize=\small' for more characters per line
\newenvironment{Shaded}{}{}
\newcommand{\AlertTok}[1]{\textcolor[rgb]{1.00,0.00,0.00}{\textbf{#1}}}
\newcommand{\AnnotationTok}[1]{\textcolor[rgb]{0.38,0.63,0.69}{\textbf{\textit{#1}}}}
\newcommand{\AttributeTok}[1]{\textcolor[rgb]{0.49,0.56,0.16}{#1}}
\newcommand{\BaseNTok}[1]{\textcolor[rgb]{0.25,0.63,0.44}{#1}}
\newcommand{\BuiltInTok}[1]{\textcolor[rgb]{0.00,0.50,0.00}{#1}}
\newcommand{\CharTok}[1]{\textcolor[rgb]{0.25,0.44,0.63}{#1}}
\newcommand{\CommentTok}[1]{\textcolor[rgb]{0.38,0.63,0.69}{\textit{#1}}}
\newcommand{\CommentVarTok}[1]{\textcolor[rgb]{0.38,0.63,0.69}{\textbf{\textit{#1}}}}
\newcommand{\ConstantTok}[1]{\textcolor[rgb]{0.53,0.00,0.00}{#1}}
\newcommand{\ControlFlowTok}[1]{\textcolor[rgb]{0.00,0.44,0.13}{\textbf{#1}}}
\newcommand{\DataTypeTok}[1]{\textcolor[rgb]{0.56,0.13,0.00}{#1}}
\newcommand{\DecValTok}[1]{\textcolor[rgb]{0.25,0.63,0.44}{#1}}
\newcommand{\DocumentationTok}[1]{\textcolor[rgb]{0.73,0.13,0.13}{\textit{#1}}}
\newcommand{\ErrorTok}[1]{\textcolor[rgb]{1.00,0.00,0.00}{\textbf{#1}}}
\newcommand{\ExtensionTok}[1]{#1}
\newcommand{\FloatTok}[1]{\textcolor[rgb]{0.25,0.63,0.44}{#1}}
\newcommand{\FunctionTok}[1]{\textcolor[rgb]{0.02,0.16,0.49}{#1}}
\newcommand{\ImportTok}[1]{\textcolor[rgb]{0.00,0.50,0.00}{\textbf{#1}}}
\newcommand{\InformationTok}[1]{\textcolor[rgb]{0.38,0.63,0.69}{\textbf{\textit{#1}}}}
\newcommand{\KeywordTok}[1]{\textcolor[rgb]{0.00,0.44,0.13}{\textbf{#1}}}
\newcommand{\NormalTok}[1]{#1}
\newcommand{\OperatorTok}[1]{\textcolor[rgb]{0.40,0.40,0.40}{#1}}
\newcommand{\OtherTok}[1]{\textcolor[rgb]{0.00,0.44,0.13}{#1}}
\newcommand{\PreprocessorTok}[1]{\textcolor[rgb]{0.74,0.48,0.00}{#1}}
\newcommand{\RegionMarkerTok}[1]{#1}
\newcommand{\SpecialCharTok}[1]{\textcolor[rgb]{0.25,0.44,0.63}{#1}}
\newcommand{\SpecialStringTok}[1]{\textcolor[rgb]{0.73,0.40,0.53}{#1}}
\newcommand{\StringTok}[1]{\textcolor[rgb]{0.25,0.44,0.63}{#1}}
\newcommand{\VariableTok}[1]{\textcolor[rgb]{0.10,0.09,0.49}{#1}}
\newcommand{\VerbatimStringTok}[1]{\textcolor[rgb]{0.25,0.44,0.63}{#1}}
\newcommand{\WarningTok}[1]{\textcolor[rgb]{0.38,0.63,0.69}{\textbf{\textit{#1}}}}
\usepackage{longtable,booktabs,array}
\newcounter{none} % for unnumbered tables
\usepackage{calc} % for calculating minipage widths
% Correct order of tables after \paragraph or \subparagraph
\usepackage{etoolbox}
\makeatletter
\patchcmd\longtable{\par}{\if@noskipsec\mbox{}\fi\par}{}{}
\makeatother
% Allow footnotes in longtable head/foot
\IfFileExists{footnotehyper.sty}{\usepackage{footnotehyper}}{\usepackage{footnote}}
\makesavenoteenv{longtable}
\setlength{\emergencystretch}{3em} % prevent overfull lines
\providecommand{\tightlist}{%
  \setlength{\itemsep}{0pt}\setlength{\parskip}{0pt}}
\usepackage{bookmark}
\IfFileExists{xurl.sty}{\usepackage{xurl}}{} % add URL line breaks if available
\urlstyle{same}
\hypersetup{
  hidelinks,
  pdfcreator={LaTeX via pandoc}}

\author{}
\date{}

\begin{document}

\textit{[Image omitted: images/053920a8b85ee3034aea7790ce00f00322054fd8a92f808c56c7fbaa3c621c7b.jpg]}

\section{usenix}\label{usenix}

THE ADVANCED

COMPUTING SYSTEMS

ASSOCIATION

\section{In Search of an Understandable Consensus
Algorithm}\label{in-search-of-an-understandable-consensus-algorithm}

Diego Ongaro and John Ousterhout, Stanford University

https://www.usenix.org/conference/atc14/technical-sessions/presentation/ongaro

\section{This paper is included in the Proceedings of USENIX ATC '14:
2014 USENIX Annual Technical
Conference.}\label{this-paper-is-included-in-the-proceedings-of-usenix-atc-14-2014-usenix-annual-technical-conference.}

June 19--20, 2014 • Philadelphia, PA

978-1-931971-10-2

Open access to the Proceedings of USENIX ATC '14: 2014 USENIX Annual
Technical Conference is sponsored by USENIX.

\section{In Search of an Understandable Consensus
Algorithm}\label{in-search-of-an-understandable-consensus-algorithm-1}

Diego Ongaro and John Ousterhout

Stanford University

\section{Abstract}\label{abstract}

Raft is a consensus algorithm for managing a replicated log. It produces
a result equivalent to (multi-)Paxos, and it is as efficient as Paxos,
but its structure is different from Paxos; this makes Raft more
understandable than Paxos and also provides a better foundation for
building practical systems. In order to enhance understandability, Raft
separates the key elements of consensus, such as leader election, log
replication, and safety, and it enforces a stronger degree of coherency
to reduce the number of states that must be considered. Results from a
user study demonstrate that Raft is easier for students to learn than
Paxos. Raft also includes a new mechanism for changing the cluster
membership, which uses overlapping majorities to guarantee safety.

\section{1 Introduction}\label{introduction}

Consensus algorithms allow a collection of machines to work as a
coherent group that can survive the failures of some of its members.
Because of this, they play a key role in building reliable large-scale
software systems. Paxos {[}13, 14{]} has dominated the discussion of
consensus algorithms over the last decade: most implementations of
consensus are based on Paxos or influenced by it, and Paxos has become
the primary vehicle used to teach students about consensus.

Unfortunately, Paxos is quite difficult to understand, in spite of
numerous attempts to make it more approachable. Furthermore, its
architecture requires complex changes to support practical systems. As a
result, both system builders and students struggle with Paxos.

After struggling with Paxos ourselves, we set out to find a new
consensus algorithm that could provide a better foundation for system
building and education. Our approach was unusual in that our primary
goal was understandability: could we define a consensus algorithm for
practical systems and describe it in a way that is significantly easier
to learn than Paxos? Furthermore, we wanted the algorithm to facilitate
the development of intuitions that are essential for system builders. It
was important not just for the algorithm to work, but for it to be
obvious why it works.

The result of this work is a consensus algorithm called Raft. In
designing Raft we applied specific techniques to improve
understandability, including decomposition (Raft separates leader
election, log replication, and safety) and state space reduction
(relative to Paxos, Raft reduces the degree of nondeterminism and the
ways servers can be inconsistent with each other). A user study with 43
students at two universities shows that Raft is significantly easier to
understand than Paxos: after learning both algorithms, 33 of these
students were able to answer questions about Raft better than questions
about Paxos.

Raft is similar in many ways to existing consensus algorithms (most
notably, Oki and Liskov's Viewstamped Replication {[}27, 20{]}), but it
has several novel features:

- Strong leader: Raft uses a stronger form of leadership than other
consensus algorithms. For example, log entries only flow from the leader
to other servers. This simplifies the management of the replicated log
and makes Raft easier to understand.

- Leader election: Raft uses randomized timers to elect leaders. This
adds only a small amount of mechanism to the heartbeats already required
for any consensus algorithm, while resolving conflicts simply and
rapidly.

- Membership changes: Raft's mechanism for changing the set of servers
in the cluster uses a new joint consensus approach where the majorities
of two different configurations overlap during transitions. This allows
the cluster to continue operating normally during configuration changes.

We believe that Raft is superior to Paxos and other consensus
algorithms, both for educational purposes and as a foundation for
implementation. It is simpler and more understandable than other
algorithms; it is described completely enough to meet the needs of a
practical system; it has several open-source implementations and is used
by several companies; its safety properties have been formally specified
and proven; and its efficiency is comparable to other algorithms.

The remainder of the paper introduces the replicated state machine
problem (Section 2), discusses the strengths and weaknesses of Paxos
(Section 3), describes our general approach to understandability
(Section 4), presents the Raft consensus algorithm (Sections 5--7),
evaluates Raft (Section 8), and discusses related work (Section 9). A
few elements of the Raft algorithm have been omitted here because of
space limitations, but they are available in an extended technical
report {[}29{]}. The additional material describes how clients interact
with the system, and how space in the Raft log can be reclaimed.

\section{2 Replicated state machines}\label{replicated-state-machines}

Consensus algorithms typically arise in the context of replicated state
machines {[}33{]}. In this approach, state machines on a collection of
servers compute identical copies of the same state and can continue
operating even if some of the servers are down. Replicated state
machines are used to solve a variety of fault tolerance problems in
distributed systems. For example, large-scale systems that

\textit{[Image omitted: images/f35d763dc25b0ccc48ffbeaad218af1ebb309c3329ead7101fd4c772872c818a.jpg]}

flowchart

\begin{Shaded}
\begin{Highlighting}[]
\NormalTok{graph TD}
\NormalTok{    A["Client"] {-}{-}\textgreater{}|①| B["Server"]}
\NormalTok{    B {-}{-}\textgreater{}|②| C["Log"]}
\NormalTok{    C {-}{-}\textgreater{}|③| D["State Machine"]}
\NormalTok{    D {-}{-}\textgreater{}|④| A}
\NormalTok{    B {-}{-}\textgreater{}|①| E["Consensus Module"]}
\NormalTok{    E {-}{-}\textgreater{}|②| C}
\NormalTok{    C {-}{-}\textgreater{}|③| D}
\NormalTok{    D {-}{-}\textgreater{}|④| A}
\NormalTok{    style A fill:\#f9f,stroke:\#333}
\NormalTok{    style B fill:\#ccf,stroke:\#333}
\NormalTok{    style C fill:\#cfc,stroke:\#333}
\NormalTok{    style D fill:\#fcc,stroke:\#333}
\NormalTok{    style E fill:\#cff,stroke:\#333}
\end{Highlighting}
\end{Shaded}

Figure 1: Replicated state machine architecture. The consensus algorithm
manages a replicated log containing state machine commands from clients.
The state machines process identical sequences of commands from the
logs, so they produce the same outputs.

have a single cluster leader, such as GFS {[}7{]}, HDFS {[}34{]}, and
RAMCloud {[}30{]}, typically use a separate replicated state machine to
manage leader election and store configuration information that must
survive leader crashes. Examples of replicated state machines include
Chubby {[}2{]} and ZooKeeper {[}9{]}.

Replicated state machines are typically implemented using a replicated
log, as shown in Figure 1. Each server stores a log containing a series
of commands, which its state machine executes in order. Each log
contains the same commands in the same order, so each state machine
processes the same sequence of commands. Since the state machines are
deterministic, each computes the same state and the same sequence of
outputs.

Keeping the replicated log consistent is the job of the consensus
algorithm. The consensus module on a server receives commands from
clients and adds them to its log. It communicates with the consensus
modules on other servers to ensure that every log eventually contains
the same requests in the same order, even if some servers fail. Once
commands are properly replicated, each server's state machine processes
them in log order, and the outputs are returned to clients. As a result,
the servers appear to form a single, highly reliable state machine.

Consensus algorithms for practical systems typically have the following
properties:

\begin{itemize}
\tightlist
\item
  They ensure safety (never returning an incorrect result) under all
  non-Byzantine conditions, including network delays, partitions, and
  packet loss, duplication, and reordering.\\
\item
  They are fully functional (available) as long as any majority of the
  servers are operational and can communicate with each other and with
  clients. Thus, a typical cluster of five servers can tolerate the
  failure of any two servers. Servers are assumed to fail by stopping;
  they may later recover from state on stable storage and rejoin the
  cluster.\\
\item
  They do not depend on timing to ensure the consistency of the logs:
  faulty clocks and extreme message delays can, at worst, cause
  availability problems.
\end{itemize}

- In the common case, a command can complete as soon as a majority of
the cluster has responded to a single round of remote procedure calls; a
minority of slow servers need not impact overall system performance.

\section{3 What's wrong with Paxos?}\label{whats-wrong-with-paxos}

Over the last ten years, Leslie Lamport's Paxos protocol {[}13{]} has
become almost synonymous with consensus: it is the protocol most
commonly taught in courses, and most implementations of consensus use it
as a starting point. Paxos first defines a protocol capable of reaching
agreement on a single decision, such as a single replicated log entry.
We refer to this subset as single-decree Paxos. Paxos then combines
multiple instances of this protocol to facilitate a series of decisions
such as a log (multi-Paxos). Paxos ensures both safety and liveness, and
it supports changes in cluster membership. Its correctness has been
proven, and it is efficient in the normal case.

Unfortunately, Paxos has two significant drawbacks. The first drawback
is that Paxos is exceptionally difficult to understand. The full
explanation {[}13{]} is notoriously opaque; few people succeed in
understanding it, and only with great effort. As a result, there have
been several attempts to explain Paxos in simpler terms {[}14, 18,
19{]}. These explanations focus on the single-decree subset, yet they
are still challenging. In an informal survey of attendees at NSDI 2012,
we found few people who were comfortable with Paxos, even among seasoned
researchers. We struggled with Paxos ourselves; we were not able to
understand the complete protocol until after reading several simplified
explanations and designing our own alternative protocol, a process that
took almost a year.

We hypothesize that Paxos' opaqueness derives from its choice of the
single-decree subset as its foundation. Single-decree Paxos is dense and
subtle: it is divided into two stages that do not have simple intuitive
explanations and cannot be understood independently. Because of this, it
is difficult to develop intuitions about why the single-decree protocol
works. The composition rules for multi-Paxos add significant additional
complexity and subtlety. We believe that the overall problem of reaching
consensus on multiple decisions (i.e., a log instead of a single entry)
can be decomposed in other ways that are more direct and obvious.

The second problem with Paxos is that it does not provide a good
foundation for building practical implementations. One reason is that
there is no widely agreed-upon algorithm for multi-Paxos. Lamport's
descriptions are mostly about single-decree Paxos; he sketched possible
approaches to multi-Paxos, but many details are missing. There have been
several attempts to flesh out and optimize Paxos, such as {[}24{]},
{[}35{]}, and {[}11{]}, but these differ from each other and from
Lamport's sketches. Systems such as Chubby {[}4{]} have implemented
Paxos-like algorithms, but in most cases their details have not been
pub-

lished.

Furthermore, the Paxos architecture is a poor one for building practical
systems; this is another consequence of the single-decree decomposition.
For example, there is little benefit to choosing a collection of log
entries independently and then melding them into a sequential log; this
just adds complexity. It is simpler and more efficient to design a
system around a log, where new entries are appended sequentially in a
constrained order. Another problem is that Paxos uses a symmetric
peer-to-peer approach at its core (though it eventually suggests a weak
form of leadership as a performance optimization). This makes sense in a
simplified world where only one decision will be made, but few practical
systems use this approach. If a series of decisions must be made, it is
simpler and faster to first elect a leader, then have the leader
coordinate the decisions.

As a result, practical systems bear little resemblance to Paxos. Each
implementation begins with Paxos, discovers the difficulties in
implementing it, and then develops a significantly different
architecture. This is time-consuming and error-prone, and the
difficulties of understanding Paxos exacerbate the problem. Paxos'
formulation may be a good one for proving theorems about its
correctness, but real implementations are so different from Paxos that
the proofs have little value. The following comment from the Chubby
implementers is typical:

There are significant gaps between the description of the Paxos
algorithm and the needs of a real-world system\ldots. the final system
will be based on an unproven protocol {[}4{]}.

Because of these problems, we concluded that Paxos does not provide a
good foundation either for system building or for education. Given the
importance of consensus in large-scale software systems, we decided to
see if we could design an alternative consensus algorithm with better
properties than Paxos. Raft is the result of that experiment.

\section{4 Designing for
understandability}\label{designing-for-understandability}

We had several goals in designing Raft: it must provide a complete and
practical foundation for system building, so that it significantly
reduces the amount of design work required of developers; it must be
safe under all conditions and available under typical operating
conditions; and it must be efficient for common operations. But our most
important goal---and most difficult challenge---was understandability.
It must be possible for a large audience to understand the algorithm
comfortably. In addition, it must be possible to develop intuitions
about the algorithm, so that system builders can make the extensions
that are inevitable in real-world implementations.

There were numerous points in the design of Raft where we had to choose
among alternative approaches. In these situations we evaluated the
alternatives based on understandability: how hard is it to explain each
alternative (for example, how complex is its state space, and does it
have subtle implications?), and how easy will it be for a reader to
completely understand the approach and its implications?

We recognize that there is a high degree of subjectivity in such
analysis; nonetheless, we used two techniques that are generally
applicable. The first technique is the well-known approach of problem
decomposition: wherever possible, we divided problems into separate
pieces that could be solved, explained, and understood relatively
independently. For example, in Raft we separated leader election, log
replication, safety, and membership changes.

Our second approach was to simplify the state space by reducing the
number of states to consider, making the system more coherent and
eliminating nondeterminism where possible. Specifically, logs are not
allowed to have holes, and Raft limits the ways in which logs can become
inconsistent with each other. Although in most cases we tried to
eliminate nondeterminism, there are some situations where nondeterminism
actually improves understandability. In particular, randomized
approaches introduce nondeterminism, but they tend to reduce the state
space by handling all possible choices in a similar fashion (``choose
any; it doesn't matter''). We used randomization to simplify the Raft
leader election algorithm.

\section{5 The Raft consensus
algorithm}\label{the-raft-consensus-algorithm}

Raft is an algorithm for managing a replicated log of the form described
in Section 2. Figure 2 summarizes the algorithm in condensed form for
reference, and Figure 3 lists key properties of the algorithm; the
elements of these figures are discussed piecewise over the rest of this
section.

Raft implements consensus by first electing a distinguished leader, then
giving the leader complete responsibility for managing the replicated
log. The leader accepts log entries from clients, replicates them on
other servers, and tells servers when it is safe to apply log entries to
their state machines. Having a leader simplifies the management of the
replicated log. For example, the leader can decide where to place new
entries in the log without consulting other servers, and data flows in a
simple fashion from the leader to other servers. A leader can fail or
become disconnected from the other servers, in which case a new leader
is elected.

Given the leader approach, Raft decomposes the consensus problem into
three relatively independent subproblems, which are discussed in the
subsections that follow:

\begin{itemize}
\tightlist
\item
  Leader election: a new leader must be chosen when an existing leader
  fails (Section 5.2).\\
\item
  Log replication: the leader must accept log entries from clients and
  replicate them across the cluster, forcing the other logs to agree
  with its own (Section 5.3).\\
\item
  Safety: the key safety property for Raft is the State Ma-
\end{itemize}

\section{State}\label{state}

\section{Persistent state on all
servers:}\label{persistent-state-on-all-servers}

(Updated on stable storage before responding to RPCs)

currentTerm

latest term server has seen (initialized to 0 on first boot, increases
monotonically)

votedFor

candidateId that received vote in current term (or null if none)

log{[}{]}

log entries; each entry contains command for state machine, and term
when entry was received by leader (first index is 1)

\section{Volatile state on all
servers:}\label{volatile-state-on-all-servers}

commitIndex

index of highest log entry known to be committed (initialized to 0,
increases monotonically)

lastApplied

index of highest log entry applied to state machine (initialized to 0,
increases monotonically)

\section{Volatile state on leaders:}\label{volatile-state-on-leaders}

(Reinitialized after election)

nextIndex{[}{]}

for each server, index of the next log entry to send to that server
(initialized to leader last log index + 1)

matchIndex{[}{]}

for each server, index of highest log entry known to be replicated on
server (initialized to 0, increases monotonically)

\section{AppendEntries RPC}\label{appendentries-rpc}

Invoked by leader to replicate log entries (§5.3); also used as
heartbeat (§5.2).

\section{Arguments:}\label{arguments}

term

leader\textquotesingle s term

leaderId

so follower can redirect clients

prevLogIndex

index of log entry immediately preceding new ones

prevLogTerm

term of prevLogIndex entry

entries{[}{]}

log entries to store (empty for heartbeat; may send more than one for
efficiency)

leaderCommit

leader\textquotesingle s commitIndex

\section{Results:}\label{results}

term

currentTerm, for leader to update itself

success

true if follower contained entry matching prevLogIndex and prevLogTerm

\section{Receiver implementation:}\label{receiver-implementation}

\begin{enumerate}
\def\labelenumi{\arabic{enumi}.}
\tightlist
\item
  Reply false if term \textless{} currentTerm (§5.1)\\
\item
  Reply false if log doesn't contain an entry at prevLogIndex whose term
  matches prevLogTerm (§5.3)\\
\item
  If an existing entry conflicts with a new one (same index but
  different terms), delete the existing entry and all that follow it
  (§5.3)\\
\item
  Append any new entries not already in the log\\
\item
  If leaderCommit \textgreater{} commitIndex, set commitIndex =
  min(leaderCommit, index of last new entry)
\end{enumerate}

\section{RequestVote RPC}\label{requestvote-rpc}

Invoked by candidates to gather votes (§5.2).

\section{Arguments:}\label{arguments-1}

term

candidate's term

candidateId

candidate requesting vote

lastLogIndex

index of candidate's last log entry (§5.4)

lastLogTerm

term of candidate's last log entry (§5.4)

\section{Results:}\label{results-1}

term

currentTerm, for candidate to update itself

voteGranted

true means candidate received vote

\section{Receiver implementation:}\label{receiver-implementation-1}

\begin{enumerate}
\def\labelenumi{\arabic{enumi}.}
\tightlist
\item
  Reply false if term \textless{} currentTerm (§5.1)\\
\item
  If votedFor is null or candidatelId, and candidate's log is at least
  as up-to-date as receiver's log, grant vote (§5.2, §5.4)
\end{enumerate}

\section{Rules for Servers}\label{rules-for-servers}

\section{All Servers:}\label{all-servers}

\begin{itemize}
\tightlist
\item
  If commitIndex \textgreater{} lastApplied: increment lastApplied,
  apply log{[}lastApplied{]} to state machine (§5.3)\\
\item
  If RPC request or response contains term T \textgreater{} currentTerm:
  set currentTerm = T, convert to follower (§5.1)
\end{itemize}

\section{Followers (§5.2):}\label{followers-5.2}

\begin{itemize}
\tightlist
\item
  Respond to RPCs from candidates and leaders\\
\item
  If election timeout elapses without receiving AppendEntries RPC from
  current leader or granting vote to candidate: convert to candidate
\end{itemize}

\section{Candidates (§5.2):}\label{candidates-5.2}

• On conversion to candidate, start election:\\
- Increment currentTerm\\
• Vote for self\\
- Reset election timer\\
- Send RequestVote RPCs to all other servers\\
• If votes received from majority of servers: become leader\\
- If AppendEntries RPC received from new leader: convert to follower\\
- If election timeout elapses: start new election

\section{Leaders:}\label{leaders}

\begin{itemize}
\tightlist
\item
  Upon election: send initial empty AppendEntries RPCs (heartbeat) to
  each server; repeat during idle periods to prevent election timeouts
  (§5.2)\\
\item
  If command received from client: append entry to local log, respond
  after entry applied to state machine (§5.3)\\
\item
  If last log index ≥ nextIndex for a follower: send AppendEntries RPC
  with log entries starting at nextIndex
\item
  If successful: update nextIndex and matchIndex for follower (§5.3)\\
\item
  If AppendEntries fails because of log inconsistency: decrement
  nextIndex and retry (§5.3)\\
\item
  If there exists an N such that N \textgreater{} commitIndex, a
  majority of matchIndex{[}i{]} ≥ N, and log{[}N{]}.term == currentTerm:
  set commitIndex = N (§5.3, §5.4).
\end{itemize}

Figure 2: A condensed summary of the Raft consensus algorithm (excluding
membership changes and log compaction). The server behavior in the
upper-left box is described as a set of rules that trigger independently
and repeatedly. Section numbers such as §5.2 indicate where particular
features are discussed. A formal specification {[}28{]} describes the
algorithm more precisely.

Election Safety: at most one leader can be elected in a given term. §5.2

Leader Append-Only: a leader never overwrites or deletes entries in its
log; it only appends new entries. §5.3

Log Matching: if two logs contain an entry with the same index and term,
then the logs are identical in all entries up through the given index.
§5.3

Leader Completeness: if a log entry is committed in a given term, then
that entry will be present in the logs of the leaders for all
higher-numbered terms. §5.4

State Machine Safety: if a server has applied a log entry at a given
index to its state machine, no other server will ever apply a different
log entry for the same index. §5.4.3

Figure 3: Raft guarantees that each of these properties is true at all
times. The section numbers indicate where each property is discussed.

chine Safety Property in Figure 3: if any server has applied a
particular log entry to its state machine, then no other server may
apply a different command for the same log index. Section 5.4 describes
how Raft ensures this property; the solution involves an additional
restriction on the election mechanism described in Section 5.2.

After presenting the consensus algorithm, this section discusses the
issue of availability and the role of timing in the system.

\section{5.1 Raft basics}\label{raft-basics}

A Raft cluster contains several servers; five is a typical number, which
allows the system to tolerate two failures. At any given time each
server is in one of three states: leader, follower, or candidate. In
normal operation there is exactly one leader and all of the other
servers are followers. Followers are passive: they issue no requests on
their own but simply respond to requests from leaders and candidates.
The leader handles all client requests (if a client contacts a follower,
the follower redirects it to the leader). The third state, candidate, is
used to elect a new leader as described in Section 5.2. Figure 4 shows
the states and their transitions; the transitions are discussed below.

Raft divides time into terms of arbitrary length, as shown in Figure 5.
Terms are numbered with consecutive integers. Each term begins with an
election, in which one or more candidates attempt to become leader as
described in Section 5.2. If a candidate wins the election, then it
serves as leader for the rest of the term. In some situations an
election will result in a split vote. In this case the term will end
with no leader; a new term (with a new election) will begin shortly.
Raft ensures that there is at most one leader in a given term.

Different servers may observe the transitions between terms at different
times, and in some situations a server may not observe an election or
even entire terms. Terms act as a logical clock {[}12{]} in Raft, and
they allow servers to detect obsolete information such as stale leaders.
Each server stores a current term number, which increases monotonically
over time. Current terms are exchanged whenever servers communicate; if
one server's current term is smaller than the other's, then it updates
its current term to the larger value. If a candidate or leader discovers
that its term is out of date, it immediately reverts to follower state.
If a server receives a request with a stale term number, it rejects the
request.

\textit{[Image omitted: images/33d6cd57880266587f5440e927ffb8e29491cbab6608ef8e89fe38a6133284e4.jpg]}

flowchart

\begin{Shaded}
\begin{Highlighting}[]
\NormalTok{graph TD}
\NormalTok{    A["Follower"] {-}{-}\textgreater{}|starts up| B["times out, starts election"]}
\NormalTok{    B {-}{-}\textgreater{}|times out, new election| C["Candidate"]}
\NormalTok{    C {-}{-}\textgreater{}|receives votes from majority of servers| D["Leader"]}
\NormalTok{    D {-}{-}\textgreater{}|discovers server with higher term| C}
\NormalTok{    C {-}{-}\textgreater{}|discovers current leader or new term| A}
\end{Highlighting}
\end{Shaded}

Figure 4: Server states. Followers only respond to requests from other
servers. If a follower receives no communication, it becomes a candidate
and initiates an election. A candidate that receives votes from a
majority of the full cluster becomes the new leader. Leaders typically
operate until they fail.

\textit{[Image omitted: images/d74968690f0e6e5d62f5e889104a48662c1ac17f76dccbf33c9a848f15206281.jpg]}

text\_image

term 1 term 2 t3 term 4 election normal operation no emerging leader
terms

Figure 5: Time is divided into terms, and each term begins with an
election. After a successful election, a single leader manages the
cluster until the end of the term. Some elections fail, in which case
the term ends without choosing a leader. The transitions between terms
may be observed at different times on different servers.

Raft servers communicate using remote procedure calls (RPCs), and the
consensus algorithm requires only two types of RPCs. RequestVote RPCs
are initiated by candidates during elections (Section 5.2), and
AppendEntries RPCs are initiated by leaders to replicate log entries and
to provide a form of heartbeat (Section 5.3). Servers retry RPCs if they
do not receive a response in a timely manner, and they issue RPCs in
parallel for best performance.

\section{5.2 Leader election}\label{leader-election}

Raft uses a heartbeat mechanism to trigger leader election. When servers
start up, they begin as followers. A server remains in follower state as
long as it receives valid RPCs from a leader or candidate. Leaders send
periodic heartbeats (AppendEntries RPCs that carry no log entries) to
all followers in order to maintain their authority. If a follower
receives no communication over a period of time called the election
timeout, then it assumes there is no viable leader and begins an
election to choose a new leader.

To begin an election, a follower increments its current term and
transitions to candidate state. It then votes for

itself and issues RequestVote RPCs in parallel to each of the other
servers in the cluster. A candidate continues in this state until one of
three things happens: (a) it wins the election, (b) another server
establishes itself as leader, or (c) a period of time goes by with no
winner. These outcomes are discussed separately in the paragraphs below.

A candidate wins an election if it receives votes from a majority of the
servers in the full cluster for the same term. Each server will vote for
at most one candidate in a given term, on a first-come-first-served
basis (note: Section 5.4 adds an additional restriction on votes). The
majority rule ensures that at most one candidate can win the election
for a particular term (the Election Safety Property in Figure 3). Once a
candidate wins an election, it becomes leader. It then sends heartbeat
messages to all of the other servers to establish its authority and
prevent new elections.

While waiting for votes, a candidate may receive an AppendEntries RPC
from another server claiming to be leader. If the leader's term
(included in its RPC) is at least as large as the candidate's current
term, then the candidate recognizes the leader as legitimate and returns
to follower state. If the term in the RPC is smaller than the
candidate's current term, then the candidate rejects the RPC and
continues in candidate state.

The third possible outcome is that a candidate neither wins nor loses
the election: if many followers become candidates at the same time,
votes could be split so that no candidate obtains a majority. When this
happens, each candidate will time out and start a new election by
incrementing its term and initiating another round of RequestVote RPCs.
However, without extra measures split votes could repeat indefinitely.

Raft uses randomized election timeouts to ensure that split votes are
rare and that they are resolved quickly. To prevent split votes in the
first place, election timeouts are chosen randomly from a fixed interval
(e.g., 150--300ms). This spreads out the servers so that in most cases
only a single server will time out; it wins the election and sends
heartbeats before any other servers time out. The same mechanism is used
to handle split votes. Each candidate restarts its randomized election
timeout at the start of an election, and it waits for that timeout to
elapse before starting the next election; this reduces the likelihood of
another split vote in the new election. Section 8.3 shows that this
approach elects a leader rapidly.

Elections are an example of how understandability guided our choice
between design alternatives. Initially we planned to use a ranking
system: each candidate was assigned a unique rank, which was used to
select between competing candidates. If a candidate discovered another
candidate with higher rank, it would return to follower state so that
the higher ranking candidate could more easily win the next election. We
found that this approach created subtle issues around availability (a
lower-ranked server might need to time out and become a candidate again
if a higher-ranked server fails, but if it does so too soon, it can
reset progress towards electing a leader). We made adjustments to the
algorithm several times, but after each adjustment new corner cases
appeared. Eventually we concluded that the randomized retry approach is
more obvious and understandable.

\textit{[Image omitted: images/510d7ec163898e40dbdf38ab1a64db12ef3218dee83f0bc2eb615fa997f2e6e5.jpg]}

bar\_stacked

{\def\LTcaptype{none} % do not increment counter
\begin{longtable}[]{@{}lll@{}}
\toprule\noalign{}
committed entries & leader & followers \\
\midrule\noalign{}
\endhead
\bottomrule\noalign{}
\endlastfoot
1 x←3 & 1 x←3 & 1 x←3 \\
2 y←1 & 1 y←1 & 1 y←1 \\
3 x←2 & 1 x←2 & 1 x←2 \\
4 x←0 & 2 x←2 & 2 x←2 \\
5 x←7 & 3 x←0 & 3 x←0 \\
6 x←5 & 3 x←5 & 3 x←5 \\
7 x←4 & 3 x←4 & 3 x←4 \\
8 x←3 & 1 x←3 & 1 x←3 \\
9 y←1 & 1 y←1 & 1 y←1 \\
10 y←9 & 1 y←9 & 1 y←9 \\
11 x←2 & 2 x←2 & 2 x←2 \\
12 x←0 & 3 x←0 & 3 x←0 \\
13 y←7 & 3 y←7 & 3 y←7 \\
14 x←5 & 3 x←5 & 3 x←5 \\
\end{longtable}
}

Figure 6: Logs are composed of entries, which are numbered sequentially.
Each entry contains the term in which it was created (the number in each
box) and a command for the state machine. An entry is considered
committed if it is safe for that entry to be applied to state machines.

\section{5.3 Log replication}\label{log-replication}

Once a leader has been elected, it begins servicing client requests.
Each client request contains a command to be executed by the replicated
state machines. The leader appends the command to its log as a new
entry, then issues AppendEntries RPCs in parallel to each of the other
servers to replicate the entry. When the entry has been safely
replicated (as described below), the leader applies the entry to its
state machine and returns the result of that execution to the client. If
followers crash or run slowly, or if network packets are lost, the
leader retries AppendEntries RPCs indefinitely (even after it has
responded to the client) until all followers eventually store all log
entries.

Logs are organized as shown in Figure 6. Each log entry stores a state
machine command along with the term number when the entry was received
by the leader. The term numbers in log entries are used to detect
inconsistencies between logs and to ensure some of the properties in
Figure 3. Each log entry also has an integer index identifying its
position in the log.

The leader decides when it is safe to apply a log entry to the state
machines; such an entry is called committed. Raft guarantees that
committed entries are durable and will eventually be executed by all of
the available state machines. A log entry is committed once the leader
that created the entry has replicated it on a majority of the servers
(e.g., entry 7 in Figure 6). This also commits

all preceding entries in the leader's log, including entries created by
previous leaders. Section 5.4 discusses some subtleties when applying
this rule after leader changes, and it also shows that this definition
of commitment is safe. The leader keeps track of the highest index it
knows to be committed, and it includes that index in future
AppendEntries RPCs (including heartbeats) so that the other servers
eventually find out. Once a follower learns that a log entry is
committed, it applies the entry to its local state machine (in log
order).

We designed the Raft log mechanism to maintain a high level of coherency
between the logs on different servers. Not only does this simplify the
system's behavior and make it more predictable, but it is an important
component of ensuring safety. Raft maintains the following properties,
which together constitute the Log Matching Property in Figure 3:

\begin{itemize}
\tightlist
\item
  If two entries in different logs have the same index and term, then
  they store the same command.\\
\item
  If two entries in different logs have the same index and term, then
  the logs are identical in all preceding entries.
\end{itemize}

The first property follows from the fact that a leader creates at most
one entry with a given log index in a given term, and log entries never
change their position in the log. The second property is guaranteed by a
simple consistency check performed by AppendEntries. When sending an
AppendEntries RPC, the leader includes the index and term of the entry
in its log that immediately precedes the new entries. If the follower
does not find an entry in its log with the same index and term, then it
refuses the new entries. The consistency check acts as an induction
step: the initial empty state of the logs satisfies the Log Matching
Property, and the consistency check preserves the Log Matching Property
whenever logs are extended. As a result, whenever AppendEntries returns
successfully, the leader knows that the follower's log is identical to
its own log up through the new entries.

During normal operation, the logs of the leader and followers stay
consistent, so the AppendEntries consistency check never fails. However,
leader crashes can leave the logs inconsistent (the old leader may not
have fully replicated all of the entries in its log). These
inconsistencies can compound over a series of leader and follower
crashes. Figure 7 illustrates the ways in which followers' logs may
differ from that of a new leader. A follower may be missing entries that
are present on the leader, it may have extra entries that are not
present on the leader, or both. Missing and extraneous entries in a log
may span multiple terms.

In Raft, the leader handles inconsistencies by forcing the followers'
logs to duplicate its own. This means that conflicting entries in
follower logs will be overwritten with entries from the leader's log.
Section 5.4 will show that this is safe when coupled with one more
restriction. To bring a follower's log into consistency with its own,
the leader must find the latest log entry where the two logs agree,
delete any entries in the follower's log after that point, and send the
follower all of the leader's entries after that point. All of these
actions happen in response to the consistency check performed by
AppendEntries RPCs. The leader maintains a nextIndex for each follower,
which is the index of the next log entry the leader will send to that
follower. When a leader first comes to power, it initializes all
nextIndex values to the index just after the last one in its log (11 in
Figure 7). If a follower's log is inconsistent with the leader's, the
AppendEntries consistency check will fail in the next AppendEntries RPC.
After a rejection, the leader decrements nextIndex and retries the
AppendEntries RPC. Eventually nextIndex will reach a point where the
leader and follower logs match. When this happens, AppendEntries will
succeed, which removes any conflicting entries in the follower's log and
appends entries from the leader's log (if any). Once AppendEntries
succeeds, the follower's log is consistent with the leader's, and it
will remain that way for the rest of the term.

\textit{[Image omitted: images/a057cb56ae5151922c170d78573decbe14121bf00ee7a7fce01d221e8f34af76.jpg]}\\
Figure 7: When the leader at the top comes to power, it is possible that
any of scenarios (a--f) could occur in follower logs. Each box
represents one log entry; the number in the box is its term. A follower
may be missing entries (a--b), may have extra uncommitted entries
(c--d), or both (e--f). For example, scenario (f) could occur if that
server was the leader for term 2, added several entries to its log, then
crashed before committing any of them; it restarted quickly, became
leader for term 3, and added a few more entries to its log; before any
of the entries in either term 2 or term 3 were committed, the server
crashed again and remained down for several terms.

The protocol can be optimized to reduce the number of rejected
AppendEntries RPCs; see {[}29{]} for details.

With this mechanism, a leader does not need to take any special actions
to restore log consistency when it comes to power. It just begins normal
operation, and the logs automatically converge in response to failures
of the Append-Entries consistency check. A leader never overwrites or
deletes entries in its own log (the Leader Append-Only Property in
Figure 3).

This log replication mechanism exhibits the desirable consensus
properties described in Section 2: Raft can accept, replicate, and apply
new log entries as long as a ma-

jority of the servers are up; in the normal case a new entry can be
replicated with a single round of RPCs to a majority of the cluster; and
a single slow follower will not impact performance.

\section{5.4 Safety}\label{safety}

The previous sections described how Raft elects leaders and replicates
log entries. However, the mechanisms described so far are not quite
sufficient to ensure that each state machine executes exactly the same
commands in the same order. For example, a follower might be unavailable
while the leader commits several log entries, then it could be elected
leader and overwrite these entries with new ones; as a result, different
state machines might execute different command sequences.

This section completes the Raft algorithm by adding a restriction on
which servers may be elected leader. The restriction ensures that the
leader for any given term contains all of the entries committed in
previous terms (the Leader Completeness Property from Figure 3). Given
the election restriction, we then make the rules for commitment more
precise. Finally, we present a proof sketch for the Leader Completeness
Property and show how it leads to correct behavior of the replicated
state machine.

\section{5.4.1 Election restriction}\label{election-restriction}

In any leader-based consensus algorithm, the leader must eventually
store all of the committed log entries. In some consensus algorithms,
such as Viewstamped Replication {[}20{]}, a leader can be elected even
if it doesn't initially contain all of the committed entries. These
algorithms contain additional mechanisms to identify the missing entries
and transmit them to the new leader, either during the election process
or shortly afterwards. Unfortunately, this results in considerable
additional mechanism and complexity. Raft uses a simpler approach where
it guarantees that all the committed entries from previous terms are
present on each new leader from the moment of its election, without the
need to transfer those entries to the leader. This means that log
entries only flow in one direction, from leaders to followers, and
leaders never overwrite existing entries in their logs.

Raft uses the voting process to prevent a candidate from winning an
election unless its log contains all committed entries. A candidate must
contact a majority of the cluster in order to be elected, which means
that every committed entry must be present in at least one of those
servers. If the candidate's log is at least as up-to-date as any other
log in that majority (where ``up-to-date'' is defined precisely below),
then it will hold all the committed entries. The RequestVote RPC
implements this restriction: the RPC includes information about the
candidate's log, and the voter denies its vote if its own log is more
up-to-date than that of the candidate.

Raft determines which of two logs is more up-to-date by comparing the
index and term of the last entries in the logs. If the logs have last
entries with different terms, then the log with the later term is more
up-to-date. If the logs end with the same term, then whichever log is
longer is more up-to-date.

\textit{[Image omitted: images/7b05c92c8e856a86b511b71dbfa15dad115a59abb0ad55b6178d315b7a973830.jpg]}

flowchart

\begin{Shaded}
\begin{Highlighting}[]
\NormalTok{graph TD}
\NormalTok{    subgraph (a)}
\NormalTok{        S1\_1["1 2"] {-}{-}\textgreater{} S1\_2["1 2"]}
\NormalTok{        S1\_2 {-}{-}\textgreater{} S2\_1["1 2"]}
\NormalTok{        S2\_1 {-}{-}\textgreater{} S3\_1["1"]}
\NormalTok{        S3\_1 {-}{-}\textgreater{} S4\_1["1"]}
\NormalTok{        S4\_1 {-}{-}\textgreater{} S5\_1["1"]}
\NormalTok{        S5\_1 {-}{-}\textgreater{} S1\_3["1 3"]}
\NormalTok{    end}
\NormalTok{    subgraph (b)}
\NormalTok{        S1\_2 {-}{-}\textgreater{} S2\_1}
\NormalTok{        S2\_1 {-}{-}\textgreater{} S3\_1}
\NormalTok{        S3\_1 {-}{-}\textgreater{} S4\_1}
\NormalTok{        S4\_1 {-}{-}\textgreater{} S5\_1}
\NormalTok{        S5\_1 {-}{-}\textgreater{} S1\_3}
\NormalTok{    end}
\NormalTok{    subgraph (c)}
\NormalTok{        S1\_3 {-}{-}\textgreater{} S2\_3["1 3"]}
\NormalTok{        S2\_3 {-}{-}\textgreater{} S3\_3["1 3"]}
\NormalTok{        S3\_3 {-}{-}\textgreater{} S4\_3["1 3"]}
\NormalTok{        S4\_3 {-}{-}\textgreater{} S5\_3["1 3"]}
\NormalTok{    end}
\NormalTok{    subgraph (d)}
\NormalTok{        S1\_3 {-}{-}\textgreater{} S2\_3}
\NormalTok{        S2\_3 {-}{-}\textgreater{} S3\_3}
\NormalTok{        S3\_3 {-}{-}\textgreater{} S4\_3}
\NormalTok{        S4\_3 {-}{-}\textgreater{} S5\_3}
\NormalTok{        S5\_3 {-}{-}\textgreater{} S1\_3}
\NormalTok{    end}
\NormalTok{    subgraph (e)}
\NormalTok{        S1\_3 {-}{-}\textgreater{} S2\_3}
\NormalTok{        S2\_3 {-}{-}\textgreater{} S3\_3}
\NormalTok{        S3\_3 {-}{-}\textgreater{} S4\_3}
\NormalTok{        S4\_3 {-}{-}\textgreater{} S5\_3}
\NormalTok{        S5\_3 {-}{-}\textgreater{} S1\_3}
\NormalTok{    end}
\end{Highlighting}
\end{Shaded}

Figure 8: A time sequence showing why a leader cannot determine
commitment using log entries from older terms. In (a) S1 is leader and
partially replicates the log entry at index 2. In (b) S1 crashes; S5 is
elected leader for term 3 with votes from S3, S4, and itself, and
accepts a different entry at log index 2. In (c) S5 crashes; S1
restarts, is elected leader, and continues replication. At this point,
the log entry from term 2 has been replicated on a majority of the
servers, but it is not committed. If S1 crashes as in (d), S5 could be
elected leader (with votes from S2, S3, and S4) and overwrite the entry
with its own entry from term 3. However, if S1 replicates an entry from
its current term on a majority of the servers before crashing, as in
(e), then this entry is committed (S5 cannot win an election). At this
point all preceding entries in the log are committed as well.

\section{5.4.2 Committing entries from previous
terms}\label{committing-entries-from-previous-terms}

As described in Section 5.3, a leader knows that an entry from its
current term is committed once that entry is stored on a majority of the
servers. If a leader crashes before committing an entry, future leaders
will attempt to finish replicating the entry. However, a leader cannot
immediately conclude that an entry from a previous term is committed
once it is stored on a majority of servers. Figure 8 illustrates a
situation where an old log entry is stored on a majority of servers, yet
can still be overwritten by a future leader.

To eliminate problems like the one in Figure 8, Raft never commits log
entries from previous terms by counting replicas. Only log entries from
the leader's current term are committed by counting replicas; once an
entry from the current term has been committed in this way, then all
prior entries are committed indirectly because of the Log Matching
Property. There are some situations where a leader could safely conclude
that an older log entry is committed (for example, if that entry is
stored on every server), but Raft takes a more conservative approach for
simplicity.

Raft incurs this extra complexity in the commitment rules because log
entries retain their original term numbers when a leader replicates
entries from previous terms. In other consensus algorithms, if a new
leader rereplicates entries from prior ``terms,'' it must do so with its
new ``term number.'' Raft's approach makes it easier

\textit{[Image omitted: images/fa7f142df0b4c80a229c7d20899f2fa1bb37235aa87b45b326e75e4fc27c8acf.jpg]}

flowchart

RPC configuration flowchart showing AE, AE, RV, and RV connections
between S1-S5 with associated RPC terms T and U

Figure 9: If S1 (leader for term T) commits a new log entry from its
term, and S5 is elected leader for a later term U, then there must be at
least one server (S3) that accepted the log entry and also voted for S5.

to reason about log entries, since they maintain the same term number
over time and across logs. In addition, new leaders in Raft send fewer
log entries from previous terms than in other algorithms (other
algorithms must send redundant log entries to renumber them before they
can be committed).

\section{5.4.3 Safety argument}\label{safety-argument}

Given the complete Raft algorithm, we can now argue more precisely that
the Leader Completeness Property holds (this argument is based on the
safety proof; see Section 8.2). We assume that the Leader Completeness
Property does not hold, then we prove a contradiction. Suppose the
leader for term T (leader \(_T\) ) commits a log entry from its term,
but that log entry is not stored by the leader of some future term.
Consider the smallest term U \textgreater{} T whose leader (leader
\(_U\) ) does not store the entry.

\begin{enumerate}
\def\labelenumi{\arabic{enumi}.}
\tightlist
\item
  The committed entry must have been absent from leader \(_{U}\) 's log
  at the time of its election (leaders never delete or overwrite
  entries).\\
\item
  leader \(_{T}\) replicated the entry on a majority of the cluster, and
  leader \(_{U}\) received votes from a majority of the cluster. Thus,
  at least one server (``the voter'') both accepted the entry from
  leader \(_{T}\) and voted for leader \(_{U}\) , as shown in Figure 9.
  The voter is key to reaching a contradiction.\\
\item
  The voter must have accepted the committed entry from leader \(_{T}\)
  before voting for leader \(_{U}\) ; otherwise it would have rejected
  the AppendEntries request from leader \(_{T}\) (its current term would
  have been higher than T).\\
\item
  The voter still stored the entry when it voted for leader \(_{U}\) ,
  since every intervening leader contained the entry (by assumption),
  leaders never remove entries, and followers only remove entries if
  they conflict with the leader.\\
\item
  The voter granted its vote to leader \(_{U}\) , so leader \(_{U}\) 's
  log must have been as up-to-date as the voter's. This leads to one of
  two contradictions.\\
\item
  First, if the voter and leader \(_{U}\) shared the same last log term,
  then leader \(_{U}\) 's log must have been at least as long as the
  voter's, so its log contained every entry in the voter's log. This is
  a contradiction, since the voter contained the committed entry and
  leader \(_{U}\) was assumed not to.\\
\item
  Otherwise, leader \(_{U}\) 's last log term must have been
\end{enumerate}

larger than the voter's. Moreover, it was larger than T, since the
voter's last log term was at least T (it contains the committed entry
from term T). The earlier leader that created leader \(_{U}\) 's last
log entry must have contained the committed entry in its log (by
assumption). Then, by the Log Matching Property, leader \(_{U}\) 's log
must also contain the committed entry, which is a contradiction.

\begin{enumerate}
\def\labelenumi{\arabic{enumi}.}
\setcounter{enumi}{7}
\tightlist
\item
  This completes the contradiction. Thus, the leaders of all terms
  greater than T must contain all entries from term T that are committed
  in term T.\\
\item
  The Log Matching Property guarantees that future leaders will also
  contain entries that are committed indirectly, such as index 2 in
  Figure 8(d).
\end{enumerate}

Given the Leader Completeness Property, it is easy to prove the State
Machine Safety Property from Figure 3 and that all state machines apply
the same log entries in the same order (see {[}29{]}).

\section{5.5 Follower and candidate
crashes}\label{follower-and-candidate-crashes}

Until this point we have focused on leader failures. Follower and
candidate crashes are much simpler to handle than leader crashes, and
they are both handled in the same way. If a follower or candidate
crashes, then future RequestVote and AppendEntries RPCs sent to it will
fail. Raft handles these failures by retrying indefinitely; if the
crashed server restarts, then the RPC will complete successfully. If a
server crashes after completing an RPC but before responding, then it
will receive the same RPC again after it restarts. Raft RPCs are
idempotent, so this causes no harm. For example, if a follower receives
an AppendEntries request that includes log entries already present in
its log, it ignores those entries in the new request.

\section{5.6 Timing and availability}\label{timing-and-availability}

One of our requirements for Raft is that safety must not depend on
timing: the system must not produce incorrect results just because some
event happens more quickly or slowly than expected. However,
availability (the ability of the system to respond to clients in a
timely manner) must inevitably depend on timing. For example, if message
exchanges take longer than the typical time between server crashes,
candidates will not stay up long enough to win an election; without a
steady leader, Raft cannot make progress.

Leader election is the aspect of Raft where timing is most critical.
Raft will be able to elect and maintain a steady leader as long as the
system satisfies the following timing requirement:

\[
b r o a d c a s t T i m e \ll e l e c t i o n T i m e o u t \ll M T B F
\]

In this inequality broadcastTime is the average time it takes a server
to send RPCs in parallel to every server in the cluster and receive
their responses; electionTimeout is the election timeout described in
Section 5.2; and MTBF is the average time between failures for a single

server. The broadcast time should be an order of magnitude less than the
election timeout so that leaders can reliably send the heartbeat
messages required to keep followers from starting elections; given the
randomized approach used for election timeouts, this inequality also
makes split votes unlikely. The election timeout should be a few orders
of magnitude less than MTBF so that the system makes steady progress.
When the leader crashes, the system will be unavailable for roughly the
election timeout; we would like this to represent only a small fraction
of overall time.

The broadcast time and MTBF are properties of the underlying system,
while the election timeout is something we must choose. Raft's RPCs
typically require the recipient to persist information to stable
storage, so the broadcast time may range from 0.5ms to 20ms, depending
on storage technology. As a result, the election timeout is likely to be
somewhere between 10ms and 500ms. Typical server MTBFs are several
months or more, which easily satisfies the timing requirement.

\section{6 Cluster membership changes}\label{cluster-membership-changes}

Up until now we have assumed that the cluster configuration (the set of
servers participating in the consensus algorithm) is fixed. In practice,
it will occasionally be necessary to change the configuration, for
example to replace servers when they fail or to change the degree of
replication. Although this can be done by taking the entire cluster
off-line, updating configuration files, and then restarting the cluster,
this would leave the cluster unavailable during the changeover. In
addition, if there are any manual steps, they risk operator error. In
order to avoid these issues, we decided to automate configuration
changes and incorporate them into the Raft consensus algorithm.

For the configuration change mechanism to be safe, there must be no
point during the transition where it is possible for two leaders to be
elected for the same term. Unfortunately, any approach where servers
switch directly from the old configuration to the new configuration is
unsafe. It isn't possible to atomically switch all of the servers at
once, so the cluster can potentially split into two independent
majorities during the transition (see Figure 10).

In order to ensure safety, configuration changes must use a two-phase
approach. There are a variety of ways to implement the two phases. For
example, some systems (e.g., {[}20{]}) use the first phase to disable
the old configuration so it cannot process client requests; then the
second phase enables the new configuration. In Raft the cluster first
switches to a transitional configuration we call joint consensus; once
the joint consensus has been committed, the system then transitions to
the new configuration. The joint consensus combines both the old and new
configurations:

- Log entries are replicated to all servers in both config-

\textit{[Image omitted: images/91869e513cbe812b41222d365bc53672b90519e7fa577793ff2d6d3f17eeab76.jpg]}

bar\_stacked

{\def\LTcaptype{none} % do not increment counter
\begin{longtable}[]{@{}lll@{}}
\toprule\noalign{}
Server & C\_old & C\_new \\
\midrule\noalign{}
\endhead
\bottomrule\noalign{}
\endlastfoot
Server 1 & 1 & 1 \\
Server 2 & 1 & 1 \\
Server 3 & 1 & 1 \\
Server 4 & 1 & 1 \\
Server 5 & 1 & 1 \\
\end{longtable}
}

Figure 10: Switching directly from one configuration to another is
unsafe because different servers will switch at different times. In this
example, the cluster grows from three servers to five. Unfortunately,
there is a point in time where two different leaders can be elected for
the same term, one with a majority of the old configuration (
\(C_{old}\) ) and another with a majority of the new configuration (
\(C_{new}\) ).

urations.

\begin{itemize}
\tightlist
\item
  Any server from either configuration may serve as leader.\\
\item
  Agreement (for elections and entry commitment) requires separate
  majorities from both the old and new configurations.
\end{itemize}

The joint consensus allows individual servers to transition between
configurations at different times without compromising safety.
Furthermore, joint consensus allows the cluster to continue servicing
client requests throughout the configuration change.

Cluster configurations are stored and communicated using special entries
in the replicated log; Figure 11 illustrates the configuration change
process. When the leader receives a request to change the configuration
from \(C_{old}\) to \(C_{new}\) , it stores the configuration for joint
consensus ( \(C_{old,new}\) in the figure) as a log entry and replicates
that entry using the mechanisms described previously. Once a given
server adds the new configuration entry to its log, it uses that
configuration for all future decisions (a server always uses the latest
configuration in its log, regardless of whether the entry is committed).
This means that the leader will use the rules of \(C_{old,new}\) to
determine when the log entry for \(C_{old,new}\) is committed. If the
leader crashes, a new leader may be chosen under either \(C_{old}\) or
\(C_{old,new}\) , depending on whether the winning candidate has
received \(C_{old,new}\) . In any case, \(C_{new}\) cannot make
unilateral decisions during this period.

Once \(C_{old,new}\) has been committed, neither \(C_{old}\) nor
\(C_{new}\) can make decisions without approval of the other, and the
Leader Completeness Property ensures that only servers with the
\(C_{old,new}\) log entry can be elected as leader. It is now safe for
the leader to create a log entry describing \(C_{new}\) and replicate it
to the cluster. Again, this configuration will take effect on each
server as soon as it is seen. When the new configuration has been
committed under the rules of \(C_{new}\) , the old configuration is
irrelevant and servers not in the new configuration can be shut down. As

\textit{[Image omitted: images/0acea3af0c577d9a1193d41c5cefd21e3325888f2b72de4066c3adeb322ad726.jpg]}

flowchart

\begin{Shaded}
\begin{Highlighting}[]
\NormalTok{graph TD}
\NormalTok{    A["C\_old can make decisions alone"] {-}{-}\textgreater{} B["C\_new can make decisions alone"]}
\NormalTok{    B {-}{-}\textgreater{} C["C\_old, new"]}
\NormalTok{    C {-}{-}\textgreater{} D["C\_old, new entry committed"]}
\NormalTok{    D {-}{-}\textgreater{} E["C\_new entry committed"]}
\NormalTok{    E {-}{-}\textgreater{} F["leader not in C\_new steps down here"]}
\NormalTok{    style E stroke:\#ff0000,stroke{-}width:2px}
\end{Highlighting}
\end{Shaded}

Figure 11: Timeline for a configuration change. Dashed lines show
configuration entries that have been created but not committed, and
solid lines show the latest committed configuration entry. The leader
first creates the \(C_{old,new}\) configuration entry in its log and
commits it to \(C_{old,new}\) (a majority of \(C_{old}\) and a majority
of \(C_{new}\) ). Then it creates the \(C_{new}\) entry and commits it
to a majority of \(C_{new}\) . There is no point in time in which
\(C_{old}\) and \(C_{new}\) can both make decisions independently.

shown in Figure 11, there is no time when \(C_{old}\) and \(C_{new}\)
can both make unilateral decisions; this guarantees safety.

There are three more issues to address for reconfiguration. The first
issue is that new servers may not initially store any log entries. If
they are added to the cluster in this state, it could take quite a while
for them to catch up, during which time it might not be possible to
commit new log entries. In order to avoid availability gaps, Raft
introduces an additional phase before the configuration change, in which
the new servers join the cluster as non-voting members (the leader
replicates log entries to them, but they are not considered for
majorities). Once the new servers have caught up with the rest of the
cluster, the reconfiguration can proceed as described above.

The second issue is that the cluster leader may not be part of the new
configuration. In this case, the leader steps down (returns to follower
state) once it has committed the \(C_{new}\) log entry. This means that
there will be a period of time (while it is committing \(C_{new}\) )
when the leader is managing a cluster that does not include itself; it
replicates log entries but does not count itself in majorities. The
leader transition occurs when \(C_{new}\) is committed because this is
the first point when the new configuration can operate independently (it
will always be possible to choose a leader from \(C_{new}\) ). Before
this point, it may be the case that only a server from \(C_{old}\) can
be elected leader.

The third issue is that removed servers (those not in \(C_{new}\) ) can
disrupt the cluster. These servers will not receive heartbeats, so they
will time out and start new elections. They will then send RequestVote
RPCs with new term numbers, and this will cause the current leader to
revert to follower state. A new leader will eventually be elected, but
the removed servers will time out again and the process will repeat,
resulting in poor availability.

To prevent this problem, servers disregard RequestVote RPCs when they
believe a current leader exists. Specifically, if a server receives a
RequestVote RPC within the minimum election timeout of hearing from a
current leader, it does not update its term or grant its vote. This does
not affect normal elections, where each server waits at least a minimum
election timeout before starting an election. However, it helps avoid
disruptions from removed servers: if a leader is able to get heartbeats
to its cluster, then it will not be deposed by larger term numbers.

\section{7 Clients and log compaction}\label{clients-and-log-compaction}

This section has been omitted due to space limitations, but the material
is available in the extended version of this paper {[}29{]}. It
describes how clients interact with Raft, including how clients find the
cluster leader and how Raft supports linearizable semantics {[}8{]}. The
extended version also describes how space in the replicated log can be
reclaimed using a snapshotting approach. These issues apply to all
consensus-based systems, and Raft's solutions are similar to other
systems.

\section{8 Implementation and
evaluation}\label{implementation-and-evaluation}

We have implemented Raft as part of a replicated state machine that
stores configuration information for RAMCloud {[}30{]} and assists in
failover of the RAMCloud coordinator. The Raft implementation contains
roughly 2000 lines of C++ code, not including tests, comments, or blank
lines. The source code is freely available {[}21{]}. There are also
about 25 independent third-party open source implementations {[}31{]} of
Raft in various stages of development, based on drafts of this paper.
Also, various companies are deploying Raft-based systems {[}31{]}.

The remainder of this section evaluates Raft using three criteria:
understandability, correctness, and performance.

\section{8.1 Understandability}\label{understandability}

To measure Raft's understandability relative to Paxos, we conducted an
experimental study using upper-level undergraduate and graduate students
in an Advanced Operating Systems course at Stanford University and a
Distributed Computing course at U.C. Berkeley. We recorded a video
lecture of Raft and another of Paxos, and created corresponding quizzes.
The Raft lecture covered the content of this paper; the Paxos lecture
covered enough material to create an equivalent replicated state
machine, including single-decree Paxos, multi-decree Paxos,
reconfiguration, and a few optimizations needed in practice (such as
leader election). The quizzes tested basic understanding of the
algorithms and also required students to reason about corner cases. Each
student watched one video, took the corresponding quiz, watched the
second video, and took the second quiz. About half of the participants
did the Paxos portion first and the other half did the Raft portion
first in order to account for both individual differences in performance
and experience gained from the first portion of the study. We compared
participants' scores on each quiz to determine whether participants
showed a better understanding of Raft.

We tried to make the comparison between Paxos and

\textit{[Image omitted: images/d86038d20e8e3868f6beb27d210910881cc5538213b9ec381d1363ef4df77373.jpg]}

scatter

{\def\LTcaptype{none} % do not increment counter
\begin{longtable}[]{@{}lll@{}}
\toprule\noalign{}
Paxos grade & Raft grade & Group \\
\midrule\noalign{}
\endhead
\bottomrule\noalign{}
\endlastfoot
10 & 42 & Raft then Paxos \\
12 & 28 & Paxos then Raft \\
15 & 22 & Paxos then Raft \\
18 & 25 & Paxos then Raft \\
20 & 30 & Paxos then Raft \\
22 & 35 & Paxos then Raft \\
25 & 38 & Paxos then Raft \\
28 & 45 & Paxos then Paxos \\
30 & 35 & Paxos then Raft \\
32 & 30 & Paxos then Raft \\
35 & 38 & Paxos then Raft \\
38 & 42 & Paxos then Raft \\
40 & 40 & Paxos then Raft \\
42 & 35 & Paxos then Raft \\
45 & 30 & Paxos then Raft \\
48 & 25 & Paxos then Raft \\
50 & 20 & Paxos then Raft \\
52 & 15 & Paxos then Raft \\
55 & 10 & Paxos then Raft \\
58 & 5 & Paxos then Raft \\
60 & 0 & Paxos then Raft \\
\end{longtable}
}

Figure 12: A scatter plot comparing 43 participants' performance on the
Raft and Paxos quizzes. Points above the diagonal (33) represent
participants who scored higher for Raft.

Raft as fair as possible. The experiment favored Paxos in two ways: 15
of the 43 participants reported having some prior experience with Paxos,
and the Paxos video is 14\% longer than the Raft video. As summarized in
Table 1, we have taken steps to mitigate potential sources of bias. All
of our materials are available for review {[}26, 28{]}.

On average, participants scored 4.9 points higher on the Raft quiz than
on the Paxos quiz (out of a possible 60 points, the mean Raft score was
25.7 and the mean Paxos score was 20.8); Figure 12 shows their
individual scores. A paired t-test states that, with 95\% confidence,
the true distribution of Raft scores has a mean at least 2.5 points
larger than the true distribution of Paxos scores.

We also created a linear regression model that predicts a new student's
quiz scores based on three factors: which quiz they took, their degree
of prior Paxos experience, and the order in which they learned the
algorithms. The model predicts that the choice of quiz produces a
12.5-point difference in favor of Raft. This is significantly higher
than the observed difference of 4.9 points, because many of the actual
students had prior Paxos experience, which helped Paxos considerably,
whereas it helped Raft slightly less. Curiously, the model also predicts
scores 6.3 points lower on Raft for people that have already taken the
Paxos quiz; although we don't know why, this does appear to be
statistically significant.

We also surveyed participants after their quizzes to see which algorithm
they felt would be easier to implement or explain; these results are
shown in Figure 13. An overwhelming majority of participants reported
Raft would be easier to implement and explain (33 of 41 for each
question). However, these self-reported feelings may be less reliable
than participants' quiz scores, and participants may have been biased by
knowledge of our hypothesis that Raft is easier to understand.

\textit{[Image omitted: images/03865faa89da05a75525ee1f31cf8f15a1bd33176796ccce4e466aab106b645c.jpg]}

bar

{\def\LTcaptype{none} % do not increment counter
\begin{longtable}[]{@{}
  >{\raggedright\arraybackslash}p{(\linewidth - 10\tabcolsep) * \real{0.0638}}
  >{\raggedright\arraybackslash}p{(\linewidth - 10\tabcolsep) * \real{0.1809}}
  >{\raggedright\arraybackslash}p{(\linewidth - 10\tabcolsep) * \real{0.2340}}
  >{\raggedright\arraybackslash}p{(\linewidth - 10\tabcolsep) * \real{0.1383}}
  >{\raggedright\arraybackslash}p{(\linewidth - 10\tabcolsep) * \real{0.2128}}
  >{\raggedright\arraybackslash}p{(\linewidth - 10\tabcolsep) * \real{0.1702}}@{}}
\toprule\noalign{}
\begin{minipage}[b]{\linewidth}\raggedright
\end{minipage} & \begin{minipage}[b]{\linewidth}\raggedright
Paxos much easier
\end{minipage} & \begin{minipage}[b]{\linewidth}\raggedright
Paxos somewhat easier
\end{minipage} & \begin{minipage}[b]{\linewidth}\raggedright
Roughly equal
\end{minipage} & \begin{minipage}[b]{\linewidth}\raggedright
Raft somewhat easier
\end{minipage} & \begin{minipage}[b]{\linewidth}\raggedright
Raft much easier
\end{minipage} \\
\midrule\noalign{}
\endhead
\bottomrule\noalign{}
\endlastfoot
implement & 1 & 4 & 3 & 20 & 13 \\
explain & 0 & 3 & 5 & 15 & 18 \\
\end{longtable}
}

Figure 13: Using a 5-point scale, participants were asked (left) which
algorithm they felt would be easier to implement in a functioning,
correct, and efficient system, and (right) which would be easier to
explain to a CS graduate student.

A detailed discussion of the Raft user study is available at {[}28{]}.

\section{8.2 Correctness}\label{correctness}

We have developed a formal specification and a proof of safety for the
consensus mechanism described in Section 5. The formal specification
{[}28{]} makes the information summarized in Figure 2 completely precise
using the TLA+ specification language {[}15{]}. It is about 400 lines
long and serves as the subject of the proof. It is also useful on its
own for anyone implementing Raft. We have mechanically proven the Log
Completeness Property using the TLA proof system {[}6{]}. However, this
proof relies on invariants that have not been mechanically checked (for
example, we have not proven the type safety of the specification).
Furthermore, we have written an informal proof {[}28{]} of the State
Machine Safety property which is complete (it relies on the
specification alone) and relatively precise (it is about 3500 words
long).

\section{8.3 Performance}\label{performance}

Raft's performance is similar to other consensus algorithms such as
Paxos. The most important case for performance is when an established
leader is replicating new log entries. Raft achieves this using the
minimal number of messages (a single round-trip from the leader to half
the cluster). It is also possible to further improve Raft's performance.
For example, it easily supports batching and pipelining requests for
higher throughput and lower latency. Various optimizations have been
proposed in the literature for other algorithms; many of these could be
applied to Raft, but we leave this to future work.

We used our Raft implementation to measure the performance of Raft's
leader election algorithm and answer two questions. First, does the
election process converge

Concern

Steps taken to mitigate bias

Materials for review {[}26, 28{]}

Equal lecture quality

Same lecturer for both. Paxos lecture based on and improved from
existing materials used in several universities. Paxos lecture is 14\%
longer.

videos

Equal quiz difficulty

Questions grouped in difficulty and paired across exams.

quizzes

Fair grading

Used rubric. Graded in random order, alternating between quizzes.

rubric

Table 1: Concerns of possible bias against Paxos in the study, steps
taken to counter each, and additional materials available.

\textit{[Image omitted: images/785fd0ca6be7b1ebb7e5f2d5c28006431a89fc9568996b4bec52e452142a161c.jpg]}\\
Figure 14: The time to detect and replace a crashed leader. The top
graph varies the amount of randomness in election timeouts, and the
bottom graph scales the minimum election timeout. Each line represents
1000 trials (except for 100 trials for ``150--150ms'') and corresponds
to a particular choice of election timeouts; for example, ``150--155ms''
means that election timeouts were chosen randomly and uniformly between
150ms and 155ms. The measurements were taken on a cluster of five
servers with a broadcast time of roughly 15ms. Results for a cluster of
nine servers are similar.

quickly? Second, what is the minimum downtime that can be achieved after
leader crashes?

To measure leader election, we repeatedly crashed the leader of a
cluster of five servers and timed how long it took to detect the crash
and elect a new leader (see Figure 14). To generate a worst-case
scenario, the servers in each trial had different log lengths, so some
candidates were not eligible to become leader. Furthermore, to encourage
split votes, our test script triggered a synchronized broadcast of
heartbeat RPCs from the leader before terminating its process (this
approximates the behavior of the leader replicating a new log entry
prior to crashing). The leader was crashed uniformly randomly within its
heartbeat interval, which was half of the minimum election timeout for
all tests. Thus, the smallest possible downtime was about half of the
minimum election timeout.

The top graph in Figure 14 shows that a small amount of randomization in
the election timeout is enough to avoid split votes in elections. In the
absence of randomness, leader election consistently took longer than 10
seconds in our tests due to many split votes. Adding just 5ms of
randomness helps significantly, resulting in a median downtime of 287ms.
Using more randomness improves worst-case behavior: with 50ms of
randomness the worst-case completion time (over 1000 trials) was 513ms.

The bottom graph in Figure 14 shows that downtime can be reduced by
reducing the election timeout. With an election timeout of 12--24ms, it
takes only 35ms on average to elect a leader (the longest trial took
152ms). However, lowering the timeouts beyond this point violates Raft's
timing requirement: leaders have difficulty broadcasting heartbeats
before other servers start new elections. This can cause unnecessary
leader changes and lower overall system availability. We recommend using
a conservative election timeout such as 150--300ms; such timeouts are
unlikely to cause unnecessary leader changes and will still provide good
availability.

\section{9 Related work}\label{related-work}

There have been numerous publications related to consensus algorithms,
many of which fall into one of the following categories:

- Lamport's original description of Paxos {[}13{]}, and attempts to
explain it more clearly {[}14, 18, 19{]}.

- Elaborations of Paxos, which fill in missing details and modify the
algorithm to provide a better foundation for implementation {[}24, 35,
11{]}.

- Systems that implement consensus algorithms, such as Chubby {[}2,
4{]}, ZooKeeper {[}9, 10{]}, and Spanner {[}5{]}. The algorithms for
Chubby and Spanner have not been published in detail, though both claim
to be based on Paxos. ZooKeeper's algorithm has been published in more
detail, but it is quite different from Paxos.

- Performance optimizations that can be applied to Paxos {[}16, 17, 3,
23, 1, 25{]}.

- Oki and Liskov's Viewstamped Replication (VR), an alternative approach
to consensus developed around the same time as Paxos. The original
description {[}27{]} was intertwined with a protocol for distributed
transactions, but the core consensus protocol has been separated in a
recent update {[}20{]}. VR uses a leader-based approach with many
similarities to Raft.

The greatest difference between Raft and Paxos is Raft's strong
leadership: Raft uses leader election as an essential part of the
consensus protocol, and it concentrates as much functionality as
possible in the leader. This approach results in a simpler algorithm
that is easier to understand. For example, in Paxos, leader election is
orthogonal to the basic consensus protocol: it serves only as a
performance optimization and is not required for achieving consensus.
However, this results in additional mechanism: Paxos includes both a
two-phase protocol for basic consensus and a separate mechanism for
leader election. In contrast, Raft incorporates leader election directly
into the consensus algorithm and uses it as the first of the two phases
of consensus. This results in less mechanism than in Paxos.

Like Raft, VR and ZooKeeper are leader-based and therefore share many of
Raft's advantages over Paxos. However, Raft has less mechanism that VR
or ZooKeeper because it minimizes the functionality in non-leaders. For
example, log entries in Raft flow in only one direction: outward from
the leader in AppendEntries RPCs. In VR

log entries flow in both directions (leaders can receive log entries
during the election process); this results in additional mechanism and
complexity. The published description of ZooKeeper also transfers log
entries both to and from the leader, but the implementation is
apparently more like Raft {[}32{]}.

Raft has fewer message types than any other algorithm for
consensus-based log replication that we are aware of. For example, VR
and ZooKeeper each define 10 different message types, while Raft has
only 4 message types (two RPC requests and their responses). Raft's
messages are a bit more dense than the other algorithms', but they are
simpler collectively. In addition, VR and ZooKeeper are described in
terms of transmitting entire logs during leader changes; additional
message types will be required to optimize these mechanisms so that they
are practical.

Several different approaches for cluster membership changes have been
proposed or implemented in other work, including Lamport's original
proposal {[}13{]}, VR {[}20{]}, and SMART {[}22{]}. We chose the joint
consensus approach for Raft because it leverages the rest of the
consensus protocol, so that very little additional mechanism is required
for membership changes. Lamport's \(\alpha\) -based approach was not an
option for Raft because it assumes consensus can be reached without a
leader. In comparison to VR and SMART, Raft's reconfiguration algorithm
has the advantage that membership changes can occur without limiting the
processing of normal requests; in contrast, VR stops all normal
processing during configuration changes, and SMART imposes an \(\alpha\)
-like limit on the number of outstanding requests. Raft's approach also
adds less mechanism than either VR or SMART.

\section{10 Conclusion}\label{conclusion}

Algorithms are often designed with correctness, efficiency, and/or
conciseness as the primary goals. Although these are all worthy goals,
we believe that understandability is just as important. None of the
other goals can be achieved until developers render the algorithm into a
practical implementation, which will inevitably deviate from and expand
upon the published form. Unless developers have a deep understanding of
the algorithm and can create intuitions about it, it will be difficult
for them to retain its desirable properties in their implementation.

In this paper we addressed the issue of distributed consensus, where a
widely accepted but impenetrable algorithm, Paxos, has challenged
students and developers for many years. We developed a new algorithm,
Raft, which we have shown to be more understandable than Paxos. We also
believe that Raft provides a better foundation for system building.
Using understandability as the primary design goal changed the way we
approached the design of Raft; as the design progressed we found
ourselves reusing a few techniques repeatedly, such as decomposing the
problem and simplifying the state space. These techniques not only
improved the understandability of Raft but also made it easier to
convince ourselves of its correctness.

\section{11 Acknowledgments}\label{acknowledgments}

The user study would not have been possible without the support of Ali
Ghodsi, David Mazières, and the students of CS 294-91 at Berkeley and CS
240 at Stanford. Scott Klemmer helped us design the user study, and
Nelson Ray advised us on statistical analysis. The Paxos slides for the
user study borrowed heavily from a slide deck originally created by
Lorenzo Alvisi. Special thanks go to David Mazières and Ezra Hoch for
finding subtle bugs in Raft. Many people provided helpful feedback on
the paper and user study materials, including Ed Bugnion, Michael Chan,
Hugues Evrard, Daniel Giffin, Arjun Gopalan, Jon Howell, Vimalkumar
Jeyakumar, Ankita Kejriwal, Aleksandar Kracun, Amit Levy, Joel Martin,
Satoshi Matsushita, Oleg Pesok, David Ramos, Robbert van Renesse, Mendel
Rosenblum, Nicolas Schiper, Deian Stefan, Andrew Stone, Ryan Stutsman,
David Terei, Stephen Yang, Matei Zaharia, 24 anonymous conference
reviewers (with duplicates), and especially our shepherd Eddie Kohler.
Werner Vogels tweeted a link to an earlier draft, which gave Raft
significant exposure. This work was supported by the Gigascale Systems
Research Center and the Multiscale Systems Center, two of six research
centers funded under the Focus Center Research Program, a Semiconductor
Research Corporation program, by STARnet, a Semiconductor Research
Corporation program sponsored by MARCO and DARPA, by the National
Science Foundation under Grant No.~0963859, and by grants from Facebook,
Google, Mellanox, NEC, NetApp, SAP, and Samsung. Diego Ongaro is
supported by The Junglee Corporation Stanford Graduate Fellowship.

\section{References}\label{references}

{[}1{]} BOLOSKY, W. J., BRADSHAW, D., HAAGENS, R. B., KUSTERS, N. P.,
AND LI, P. Paxos replicated state machines as the basis of a
high-performance data store. In Proc. NSDI'11, USENIX Conference on
Networked Systems Design and Implementation (2011), USENIX,
pp.~141--154.\\
{[}2{]} BURROWS, M. The Chubby lock service for loosely-coupled
distributed systems. In Proc. OSDI'06, Symposium on Operating Systems
Design and Implementation (2006), USENIX, pp.~335--350.\\
{[}3{]} CAMARGOS, L. J., SCHMIDT, R. M., AND PEDONE, F. Multicoordinated
Paxos. In Proc. PODC'07, ACM Symposium on Principles of Distributed
Computing (2007), ACM, pp.~316--317.\\
{[}4{]} CHANDRA, T. D., GRIESEMER, R., AND REDSTONE, J. Paxos made live:
an engineering perspective. In Proc. PODC'07, ACM Symposium on
Principles of Distributed Computing (2007), ACM, pp.~398--407.

{[}5{]} CORBETT, J. C., DEAN, J., EPSTEIN, M., FIKES, A., FROST, C.,
FURMAN, J. J., GHEMAWAT, S., GUBAREV, A., HEISER, C., HOCHSCHILD, P.,
HSIEH, W., KANTHAK, S., KOGAN, E., LI, H., LLOYD, A., MELNIK, S.,
MWAURA, D., NAGLE, D., QUINLAN, S., RAO, R., ROLIG, L., SAITO, Y.,
SZYMANIAK, M., TAYLOR, C., WANG, R., AND WOODFORD, D. Spanner: Google's
globally-distributed database. In Proc. OSDI'12, USENIX Conference on
Operating Systems Design and Implementation (2012), USENIX,
pp.~251--264.\\
{[}6{]} COUSINEAU, D., DOLIGEZ, D., LAMPORT, L., MERZ, S., RICKETTS, D.,
AND VANZETTO, H. TLA \(^{+}\) proofs. In Proc. FM'12, Symposium on
Formal Methods (2012), D. Giannakopoulou and D. Méry, Eds., vol.~7436 of
Lecture Notes in Computer Science, Springer, pp.~147--154.\\
{[}7{]} GHEMAWAT, S., GOBIOFF, H., AND LEUNG, S.-T. The Google file
system. In Proc. SOSP'03, ACM Symposium on Operating Systems Principles
(2003), ACM, pp.~29--43.\\
{[}8{]} HERLIHY, M. P., AND WING, J. M. Linearizability: a correctness
condition for concurrent objects. ACM Transactions on Programming
Languages and Systems 12 (July 1990), 463--492.\\
{[}9{]} HUNT, P., KONAR, M., JUNQUEIRA, F. P., AND REED, B. ZooKeeper:
wait-free coordination for internet-scale systems. In Proc ATC'10,
USENIX Annual Technical Conference (2010), USENIX, pp.~145--158.\\
{[}10{]} JUNQUEIRA, F. P., REED, B. C., AND SERAFINI, M. Zab:
High-performance broadcast for primary-backup systems. In Proc. DSN'11,
IEEE/IFIP Int'l Conf. on Dependable Systems \& Networks (2011), IEEE
Computer Society, pp.~245--256.\\
{[}11{]} KIRSCH, J., AND AMIR, Y. Paxos for system builders. Tech.
Rep.~CNDS-2008-2, Johns Hopkins University, 2008.\\
{[}12{]} LAMPORT, L. Time, clocks, and the ordering of events in a
distributed system. Communications of the ACM 21, 7 (July 1978),
558--565.\\
{[}13{]} LAMPORT, L. The part-time parliament. ACM Transactions on
Computer Systems 16, 2 (May 1998), 133--169.\\
{[}14{]} LAMPORT, L. Paxos made simple. ACM SIGACT News 32, 4
(Dec.~2001), 18--25.\\
{[}15{]} LAMPORT, L. Specifying Systems, The TLA+ Language and Tools for
Hardware and Software Engineers. Addison-Wesley, 2002.\\
{[}16{]} LAMPORT, L. Generalized consensus and Paxos. Tech.
Rep.~MSR-TR-2005-33, Microsoft Research, 2005.\\
{[}17{]} LAMPORT, L. Fast paxos. Distributed Computing 19, 2 (2006),
79--103.\\
{[}18{]} LAMPSON, B. W. How to build a highly available system using
consensus. In Distributed Algorithms, O. Baboaglu and K. Marzullo, Eds.
Springer-Verlag, 1996, pp.~1--17.\\
{[}19{]} LAMPSON, B. W. The ABCD's of Paxos. In Proc. PODC'01, ACM
Symposium on Principles of Distributed Computing (2001), ACM,
pp.~13--13.

{[}20{]} LISKOV, B., AND COWLING, J. Viewstamped replication revisited.
Tech. Rep.~MIT-CSAIL-TR-2012-021, MIT, July 2012.\\
{[}21{]} LogCabin source code. http://github.com/logcabin/logcabin.\\
{[}22{]} LORCH, J. R., ADYA, A., BOLOSKY, W. J., CHAIKEN, R., DOUCEUR,
J. R., AND HOWELL, J. The SMART way to migrate replicated stateful
services. In Proc. EuroSys'06, ACM SIGOPS/EuroSys European Conference on
Computer Systems (2006), ACM, pp.~103--115.\\
{[}23{]} MAO, Y., JUNQUEIRA, F. P., AND MARZULLO, K. Mencius: building
efficient replicated state machines for WANs. In Proc. OSDI'08, USENIX
Conference on Operating Systems Design and Implementation (2008),
USENIX, pp.~369--384.\\
{[}24{]} MAZIÈRES, D. Paxos made practical.
http://www.scs.stanford.edu/\textasciitilde dm/home/papers/paxos.pdf,
Jan.~2007.\\
{[}25{]} MORARU, I., ANDERSEN, D. G., AND KAMINSKY, M. There is more
consensus in egalitarian parliaments. In Proc. SOSP'13, ACM Symposium on
Operating System Principles (2013), ACM.\\
{[}26{]} Raft user study.
http://ramcloud.stanford.edu/\textasciitilde ongaro/userstudy/.\\
{[}27{]} OKI, B. M., AND LISKOV, B. H. Viewstamped replication: A new
primary copy method to support highly-available distributed systems. In
Proc. PODC'88, ACM Symposium on Principles of Distributed Computing
(1988), ACM, pp.~8--17.\\
{[}28{]} ONGARO, D. Consensus: Bridging Theory and Practice. PhD thesis,
Stanford University, 2014 (work in progress).
http://ramcloud.stanford.edu/\textasciitilde ongaro/thesis.pdf.\\
{[}29{]} ONGARO, D., AND OUSTERHOUT, J. In search of an understandable
consensus algorithm (extended version).
http://ramcloud.stanford.edu/raft.pdf.\\
{[}30{]} OUSTERHOUT, J., AGRAWAL, P., ERICKSON, D., KOZYRAKIS, C.,
LEVERICH, J., MAZIÈRES, D., MITRA, S., NARAYANAN, A., ONGARO, D.,
PARULKAR, G., ROSENBLUM, M., RUMBLE, S. M., STRATMANN, E., AND STUTSMAN,
R. The case for RAMCloud. Communications of the ACM 54 (July 2011),
121--130.\\
{[}31{]} Raft consensus algorithm website.
http://raftconsensus.github.io.\\
{[}32{]} REED, B. Personal communications, May 17, 2013.\\
{[}33{]} SCHNEIDER, F. B. Implementing fault-tolerant services using the
state machine approach: a tutorial. ACM Computing Surveys 22, 4
(Dec.~1990), 299--319.\\
{[}34{]} SHVACHKO, K., KUANG, H., RADIA, S., AND CHANSLER, R. The Hadoop
distributed file system. In Proc. MSST'10, Symposium on Mass Storage
Systems and Technologies (2010), IEEE Computer Society, pp.~1--10.\\
{[}35{]} VAN RENESSE, R. Paxos made moderately complex. Tech. rep.,
Cornell University, 2012.

\end{document}
